\documentclass{situationist-preprint}

\title{Intending to Try}
\msid{ACT-11-a}
\msfamily{ACT-11}
\msdiscipline{Action, Agency and Free Will}
\mstype{Principal manuscript}
\mskeywords{intention; trying; impossibility; belief constraint; wish; planning}

\begin{document}
\maketitle

\begin{abstract}
Can one intend to do what one believes impossible? Aristotle held that wish may be for impossible things but choice may not, and the planning theory of intention gives the reason: an intention settles a practical question within the world the agent takes herself to inhabit, and an agent who believes she will not do something has nothing to settle. We argue that one cannot intend to do what one believes one will not do, and a fortiori what one believes impossible. The constraint is constitutive of intention, not merely a norm of rationality. But one can intend to try, since trying is something one may believe one can do even when one believes one will fail; and many intentions concerning uncertain outcomes, such as the athlete's or the escapee's, have this form. We show that the account explains the apparent counterexamples, including the agent who sets out to do what she thinks almost certainly beyond her, and that it agrees with the restriction of agglomeration to single plans that ACT-15 uses.
\end{abstract}

\section{The problem}

\citet[III.2, 1111b20--30]{AristotleNE} distinguishes choice from wish: one may wish for impossible things, such as immortality, but no one chooses them. The modern version of the question concerns intention. \citet{Grice1971} argued that to intend to $A$ one must believe that one will $A$; \citet{Bratman1987} replaced this with the weaker requirement that one not believe that one will not $A$, and argued that, where the agent doubts her success, she may intend to try rather than intend to succeed. \citet{McCann1986} defended the possibility of intending an outcome one believes unlikely. We argue for a definite answer: one cannot intend what one believes impossible, but one can intend to try, and many intentions are of this form.

\section{Framework}

\begin{definition}[Intention]\label{d:int}
An \emph{intention} to $A$ is an attitude whose content is a plan that includes $A$-ing and which involves commitment: it settles the question whether to $A$ and constrains further reasoning (ACT-06-a).
\end{definition}

\begin{definition}[Trying]\label{d:try}
To \emph{try} to $A$ is to do what one believes is a way of $A$-ing, or of making it more likely that one $A$s, with the aim of $A$-ing \citep{OShaughnessy1973,Hornsby1980}.
\end{definition}

\begin{principle}[Settling within belief]\label{p:settle}
An attitude settles the question whether to $A$ only if, in the world as the agent believes it to be, there is a course of conduct she may adopt that includes $A$-ing.
\end{principle}

Principle~\ref{p:settle} states what settling is. To settle whether to $A$ is to fix one's course in the world one is about to act in, and that is, from the agent's point of view, the world as she believes it to be. If, by her lights, every course open to her excludes $A$-ing, there is no course for the attitude to fix.

\section{The argument}

\begin{theorem}\label{t:main}
No agent intends to $A$ while believing that she will not $A$. In particular, no agent intends to do what she believes impossible.
\end{theorem}

\begin{proof}
Suppose an agent believes she will not $A$. Then in the world as she believes it to be, no course of conduct she adopts includes $A$-ing. By Principle~\ref{p:settle}, no attitude of hers settles whether to $A$, and by Definition~\ref{d:int}, she does not intend to $A$. If she believes $A$-ing impossible, she believes she will not $A$, and the first part applies.
\end{proof}

The attitude such an agent can have toward $A$ is Aristotle's wish: she may want to $A$, regret that she cannot, and imagine doing it. What she cannot do is make $A$-ing part of her plan. The constraint is constitutive: it is not that an agent who intends what she believes impossible is irrational, but that what she has is not an intention. It follows, a fortiori, that one cannot \emph{rationally} intend what one believes impossible, which is the form in which ACT-15-a uses the constraint.

\begin{theorem}\label{t:try}
An agent may intend to try to $A$ while believing that she will not $A$, provided she believes that she can try.
\end{theorem}

\begin{proof}
Trying to $A$ is doing what one believes is a way of $A$-ing, or of raising its probability, with the aim of $A$-ing (Definition~\ref{d:try}). An agent may believe that there is such a thing she can do even when she believes that it will not succeed: the weightlifter who believes the bar too heavy can still grip it and heave. In the world as she believes it to be, there is a course of conduct that includes trying, so Principle~\ref{p:settle} is satisfied, and nothing prevents an attitude with commitment toward that course.
\end{proof}

\paragraph{Many intentions have this form.} Where the agent is uncertain of success, without believing she will fail, she may intend either to $A$ or to try to $A$, and what she intends shows in her plans. The basketball player who shoots from half court as the clock runs out does not plan on the basket: she does not plan the celebration, and her teammates do not set up for the next play on the assumption that it went in. She intends to make the shot in the sense of trying to make it. The prisoner who plans an escape he thinks unlikely to succeed makes contingency plans for failure, and these show that his plan is built around an attempt. The content of an intention is the part of the agent's conduct she treats as settled, and where success is not treated as settled, the intention is to try. This includes the intentions of anyone who undertakes a long shot, a rescue in doubtful conditions or a difficult examination, and it is why \citet{Bratman1987} could treat the video-game player, who believes he will not hit both targets, as trying to hit each.

\section{Objections and replies}

\begin{objection}[Defiant intention]
A climber says: \enquote{I know I cannot reach the summit before dark, but I intend to.} She is not confused, and she sets out.
\end{objection}

\begin{reply}
Either she does not really believe she cannot, in which case her avowal expresses doubt rather than belief and Theorem~\ref{t:main} does not apply, or she believes it and what she intends is to try, to get as high as she can before dark. Her plans settle the question: if she carries bivouac equipment for the case in which she does not make it, her plan is built around an attempt. The avowal is idiomatic compression of an intention to try.
\end{reply}

\begin{objection}[Unrecognized impossibility]
A geometer intends to trisect an arbitrary angle with ruler and compass, not knowing that it is impossible. She intends the impossible.
\end{objection}

\begin{reply}
She has an intention that cannot be fulfilled, which is a case of error, not of intending the impossible in the sense of the question. The question concerns whether an agent can settle on what she takes to be impossible, and Theorem~\ref{t:main} answers it. Once she learns the proof, her intention lapses, or becomes an intention to try for some other purpose, such as finding where a purported construction goes wrong; this is what Theorem~\ref{t:main} predicts.
\end{reply}

\begin{objection}[Incompatible intentions]
ACT-15-a holds that an agent may intend each of two outcomes she knows cannot both occur. Then she intends, of the pair, something she believes she will not achieve.
\end{objection}

\begin{reply}
She does not intend the pair. Each of her intentions is consistent with her beliefs, since she does not believe she will fail to bring about either outcome singly, and the intentions do not belong to a single plan, so agglomeration does not require her to intend their conjunction (ACT-06-a). Theorem~\ref{t:main} constrains each intention, not the conjunction of intentions that belong to different plans.
\end{reply}

\section{Conclusion}

An intention settles what to do within the world one believes one is acting in, and so cannot be directed at what one believes one will not or cannot do. What remains toward such an object is wish, or an intention to try. Since many of our intentions concern outcomes we doubt, many of them are intentions to try. A fuller analysis of trying is given in ACT-11-b.

\begin{thebibliography}{9}\small
\bibitem[Aristotle(NE)]{AristotleNE} Aristotle. \emph{Nicomachean Ethics}.
\bibitem[Bratman(1987)]{Bratman1987} Bratman, M.~E. (1987). \emph{Intention, Plans, and Practical Reason}. Harvard University Press.
\bibitem[Grice(1971)]{Grice1971} Grice, H.~P. (1971). Intention and uncertainty. \emph{Proceedings of the British Academy}, 57, 263--279.
\bibitem[Hornsby(1980)]{Hornsby1980} Hornsby, J. (1980). \emph{Actions}. Routledge \& Kegan Paul.
\bibitem[McCann(1986)]{McCann1986} McCann, H.~J. (1986). Rationality and the range of intention. \emph{Midwest Studies in Philosophy}, 10, 191--211.
\bibitem[O'Shaughnessy(1973)]{OShaughnessy1973} O'Shaughnessy, B. (1973). Trying (as the mental \enquote{pineal gland}). \emph{The Journal of Philosophy}, 70(13), 365--386.
\end{thebibliography}

\end{document}
